\documentclass[11pt,a4paper]{article}

% ============================================================
%  PRÉAMBULE — AMBIANCE NOIRE
% ============================================================
\usepackage[utf8]{inputenc}
\usepackage[T1]{fontenc}
\usepackage[french]{babel}
\usepackage{lmodern}
\usepackage{amsmath,amssymb}
\usepackage{geometry}
\geometry{margin=2.2cm}
\usepackage{xcolor}
\usepackage[most]{tcolorbox}
\usepackage{enumitem}
\usepackage{tikz}
\usepackage{pgfplots}
\pgfplotsset{compat=1.18}
\usepackage{fancyhdr}
\usepackage{hyperref}
\usepackage{titlesec}
\usepackage{epigraph}

% ---- PALETTE : UNIQUEMENT DES DÉGRADÉS DE GRIS ----
\definecolor{ink}{HTML}{0D0D0D}          % presque noir
\definecolor{charcoal}{HTML}{262626}     % charbon
\definecolor{slate}{HTML}{4D4D4D}        % ardoise
\definecolor{gray}{HTML}{808080}         % gris moyen
\definecolor{silver}{HTML}{B3B3B3}       % argent
\definecolor{smoke}{HTML}{D9D9D9}        % fumée
\definecolor{paper}{HTML}{F2F2F2}        % papier
\definecolor{white}{HTML}{FFFFFF}        % blanc

\hypersetup{colorlinks=true,linkcolor=ink,urlcolor=slate,linktoc=all}

% ============================================================
%  ENVIRONNEMENTS — DÉGRADÉS DE GRIS
% ============================================================

% Indice : du blanc vers la fumée
\newtcolorbox{indice}[1][Indice]{
  enhanced, breakable,
  colframe=slate, coltitle=white, colbacktitle=slate,
  interior style={top color=white, bottom color=smoke, shading angle=90},
  boxrule=0.8pt, arc=2mm,
  fonttitle=\bfseries, title={#1}}

% Méthode : du papier vers l'argent
\newtcolorbox{methode}[1][Méthode d'enquête]{
  enhanced, breakable,
  colframe=charcoal, coltitle=white, colbacktitle=charcoal,
  interior style={top color=paper, bottom color=silver!50, shading angle=90},
  boxrule=0.8pt, arc=2mm,
  fonttitle=\bfseries, title={#1}}

% Rapport d'observation : du blanc vers le papier
\newtcolorbox{rapport}[1][Rapport d'observation]{
  enhanced, breakable,
  colframe=gray, coltitle=white, colbacktitle=gray,
  interior style={top color=white, bottom color=paper, shading angle=90},
  boxrule=0.8pt, arc=2mm,
  fonttitle=\bfseries, title={#1}}

% Alerte : de la fumée vers le papier (bord sombre)
\newtcolorbox{alerte}[1][Alerte]{
  enhanced, breakable,
  colframe=ink, coltitle=white, colbacktitle=ink,
  interior style={top color=smoke, bottom color=paper, shading angle=90},
  boxrule=1pt, arc=2mm,
  fonttitle=\bfseries, title={#1}}

% Affaire : du papier vers la fumée
\newtcolorbox{affaire}[1][Affaire à résoudre]{
  enhanced, breakable,
  colframe=charcoal, coltitle=white, colbacktitle=charcoal,
  interior style={top color=paper, bottom color=smoke, shading angle=90},
  boxrule=0.8pt, arc=2mm,
  fonttitle=\bfseries, title={#1}}

% Révélation : de l'argent vers la fumée
\newtcolorbox{revelation}[1][Révélation]{
  enhanced, breakable,
  colframe=ink, coltitle=white, colbacktitle=ink,
  interior style={top color=silver!70, bottom color=smoke, shading angle=90},
  boxrule=0.8pt, arc=2mm,
  fonttitle=\bfseries, title={#1}}

% ---- EN-TÊTES DE CHAPITRE (bandeau noir dégradé) ----
\newcommand{\chapterheader}[2]{%
  \clearpage
  \begin{tcolorbox}[
    enhanced,
    colback=ink, colframe=ink,
    interior style={top color=ink, bottom color=charcoal, shading angle=90},
    boxrule=0pt, arc=0pt,
    left=18pt, right=18pt, top=14pt, bottom=14pt,
    width=\textwidth,
  ]
    {\color{white}\Huge\bfseries #1}\\[6pt]
    {\color{silver}\large\itshape #2}
  \end{tcolorbox}
  \vspace{6pt}
}

% ---- EN-TÊTES / PIEDS DE PAGE ----
\pagestyle{fancy}
\fancyhf{}
\fancyhead[L]{\small\textcolor{ink}{\textbf{Dossier « Second Degré »}}}
\fancyhead[R]{\small\textcolor{slate}{Inspectrice M. Vasseur}}
\fancyfoot[C]{\textcolor{slate}{\thepage}}
\renewcommand{\headrulewidth}{0.4pt}
\renewcommand{\headrule}{\hbox to\headwidth{\color{slate}\leaders\hrule height \headrulewidth\hfill}}

\setlength{\parindent}{0pt}
\setlength{\parskip}{4pt}

% ============================================================
\begin{document}

% ============================================================
%  PAGE DE TITRE — AFFICHE NOIRE
% ============================================================
\begin{titlepage}
\pagecolor{ink}
\color{white}
\centering
\vspace*{1.5cm}

\begin{tcolorbox}[
  enhanced, colback=ink, colframe=ink,
  interior style={top color=ink, bottom color=charcoal, shading angle=90},
  boxrule=0pt, arc=0pt, width=\textwidth,
  left=20pt, right=20pt, top=20pt, bottom=20pt,
]
\centering
{\fontsize{42}{48}\selectfont\bfseries L'ODYSSÉE\\[4pt] DU SECOND DEGRÉ}\\[14pt]
{\color{silver}\Large\itshape Une enquête en six tableaux}\\[6pt]
{\color{gray}\large Inspectrice Maëlle Vasseur}
\end{tcolorbox}

\vspace{1cm}

\begin{tikzpicture}
\begin{axis}[
  width=11cm, height=6cm,
  axis lines=middle,
  xlabel={\color{white}$x$}, ylabel={\color{white}$y$},
  xmin=-3, xmax=6, ymin=-3, ymax=6,
  xtick={-2,...,5}, ytick={-2,...,5},
  grid=both, grid style={gray!30},
  axis line style={white}, tick style={white},
  xticklabel style={color=white}, yticklabel style={color=white},
]
\addplot[domain=-1.8:5.8, samples=150, very thick, white] {0.4*(x-2)^2 - 1};
\addplot[only marks, mark=*, mark size=3pt, silver] coordinates {(2,-1)};
\node at (axis cs:2.6,0.3) {\color{silver}\small $S(\alpha,\beta)$};
\end{axis}
\end{tikzpicture}

\vfill

{\color{silver}\large\itshape « Chaque parabole cache un crime.\\[2pt]
Chaque sommet cache un coupable. »}\\[0.6cm]

{\color{gray}\small D'après l'application interactive \emph{L'Odyssée du Second Degré}}\\[2pt]
{\color{gray}\small Yann Merdy — Licence non commerciale}

\end{titlepage}

\nopagecolor
\color{ink}

% ---- SOMMAIRE ----
\clearpage
\renewcommand{\contentsname}{\textcolor{ink}{Table des affaires}}
\tableofcontents
\thispagestyle{empty}

% ============================================================
%  PROLOGUE
% ============================================================
\chapterheader{Prologue}{L'appel de 23h47}

La pluie frappait la vitre du bureau 407 comme un métronome fatigué. L'Inspectrice Maëlle Vasseur n'avait pas dormi depuis trois jours. Sur la table, six dossiers. Six affaires non résolues. Six paraboles qui n'avaient jamais livré leurs secrets.

Le téléphone sonna.

— Inspectrice ? Une voix à peine audible. Six crimes. Six scènes. Une seule signature : $$f(x) = ax^2 + bx + c$$. Vous avez jusqu'à l'aube.

La communication coupa. L'Inspectrice ouvrit le premier dossier. Elle savait qu'une nuit entière ne suffirait pas — mais elle savait aussi qu'elle n'avait pas le choix.

\begin{indice}[Ce que vous devez savoir avant de commencer]
Ce carnet est le \textbf{journal de bord} de l'Inspectrice Maëlle Vasseur. Chaque chapitre est une \textbf{scène de crime} mathématique. Pour résoudre l'enquête, vous devrez maîtriser :
\begin{itemize}[leftmargin=*]
  \item la \textbf{fonction du second degré} et sa courbe, la parabole ;
  \item le \textbf{sommet}, l'\textbf{axe de symétrie}, la \textbf{forme canonique} ;
  \item les \textbf{conditions} imposées à une fonction (trajectoire, points, valeurs) ;
  \item la \textbf{rédaction rigoureuse} d'un raisonnement.
\end{itemize}
Un seul mot d'ordre : \emph{observer, déduire, conclure}.
\end{indice}

\vfill
\begin{center}
\textcolor{slate}{\itshape\small — Le dossier est ouvert. —}
\end{center}

% ============================================================
%  CHAPITRE 1 — LA FONCTION DU SECOND DEGRÉ
% ============================================================
\chapterheader{Chapitre 1}{Le portrait du suspect}

Le premier dossier s'ouvrit comme une porte qu'on n'aurait pas dû franchir. À l'intérieur, une photo : une courbe en forme de sourire, ou de moue, selon l'angle. Et sous la photo, une formule gravée comme sur une pierre tombale :
\[
  f(x) = ax^2 + bx + c, \qquad a \neq 0.
\]
L'Inspectrice reposa sa tasse. C'était le portrait du suspect. Pas n'importe quel suspect : le suspect numéro un de toutes ses affaires non résolues.

\section{Définition et forme développée}

\begin{indice}[Définition — le portrait-robot]
On appelle \textbf{fonction polynôme du second degré} toute fonction $f$ définie sur $\mathbb{R}$ par une expression de la forme
\[
  \boxed{\;f(x) = ax^2 + bx + c\;}
\]
où $a$, $b$, $c$ sont trois réels, avec la condition essentielle $a \neq 0$.
\begin{itemize}[leftmargin=*]
  \item $a$ : coefficient dominant (coefficient de $x^2$) ;
  \item $b$ : coefficient de $x$ ;
  \item $c$ : terme constant.
\end{itemize}
Cette écriture s'appelle la \textbf{forme développée} de $f$.
\end{indice}

\begin{alerte}[Piège n°1 — le faux suspect]
Si $a = 0$, la fonction devient $f(x) = bx + c$ : ce n'est plus du second degré, mais une simple droite. Le suspect s'est alors \emph{déguisé}. La condition $a \neq 0$ est donc \textbf{indispensable}.
\end{alerte}

\section{La parabole — la signature du suspect}

\begin{indice}[Représentation graphique]
La courbe représentative d'une fonction du second degré s'appelle une \textbf{parabole}. Elle possède :
\begin{itemize}[leftmargin=*]
  \item un \textbf{sommet} $S$ : le point le plus bas ou le plus haut ;
  \item un \textbf{axe de symétrie} vertical passant par le sommet.
\end{itemize}
\end{indice}

\begin{center}
\begin{tikzpicture}
\begin{axis}[
  width=13cm, height=6cm,
  axis lines=middle, xlabel=$x$, ylabel=$y$,
  xmin=-4, xmax=4, ymin=-2, ymax=8,
  xtick={-3,...,3}, ytick={-1,...,7},
  grid=both, grid style={gray!20},
  legend style={at={(0.02,0.98)},anchor=north west,font=\small},
]
\addplot[domain=-3.2:3.2, samples=150, very thick, charcoal] {x^2 - 2};
\addlegendentry{$a>0$ : vers le haut}
\addplot[domain=-3.2:3.2, samples=150, very thick, gray, dashed] {-x^2 + 3};
\addlegendentry{$a<0$ : vers le bas}
\end{axis}
\end{tikzpicture}
\end{center}

\section{Le rôle du signe de \texorpdfstring{$a$}{a}}

\begin{indice}[Orientation de la parabole]
Le signe de $a$ détermine \emph{complètement} l'orientation :
\begin{itemize}[leftmargin=*]
  \item $a > 0$ : parabole tournée \textbf{vers le haut}. Le sommet est un \emph{minimum}.
  \item $a < 0$ : parabole tournée \textbf{vers le bas}. Le sommet est un \emph{maximum}.
\end{itemize}
Plus $|a|$ est grand, plus la parabole est « resserrée ».
\end{indice}

\begin{alerte}[Piège n°2 — confusion]
Le signe de $a$ donne l'\textbf{orientation}, \textbf{pas} la position du sommet ni le signe de $f(x)$ sur tout $\mathbb{R}$.
\end{alerte}

\section{Le sommet et l'axe de symétrie}

\begin{indice}[Formules clés]
Pour $f(x) = ax^2+bx+c$ avec $a \neq 0$ :
\[
  \alpha = -\frac{b}{2a} \quad \text{(abscisse du sommet)},
  \qquad
  \beta = f(\alpha) \quad \text{(ordonnée du sommet)}.
\]
Le sommet est $S(\alpha\,;\,\beta)$. L'\textbf{axe de symétrie} est la droite
\[
  \boxed{\;x = \alpha = -\frac{b}{2a}\;}
\]
\end{indice}

\begin{methode}[Le réflexe de l'Inspectrice]
Dès qu'un énoncé parle de sommet, d'axe de symétrie ou de maximum/minimum, l'Inspectrice note immédiatement $\alpha = -\dfrac{b}{2a}$. C'est sa boussole.
\end{methode}

\section{La forme canonique — le vrai visage}

\begin{indice}[Forme canonique]
Toute fonction du second degré peut s'écrire sous la \textbf{forme canonique} :
\[
  \boxed{\;f(x) = a(x-\alpha)^2 + \beta\;}
\]
où $(\alpha\,;\,\beta)$ sont les coordonnées du sommet.
Cette forme met en évidence :
\begin{itemize}[leftmargin=*]
  \item le \textbf{sommet} $S(\alpha\,;\,\beta)$ ;
  \item l'\textbf{axe de symétrie} $x = \alpha$ ;
  \item l'\textbf{orientation} donnée par le signe de $a$.
\end{itemize}
\end{indice}

\begin{alerte}[Piège n°3 — la fausse identité]
La forme canonique n'est \textbf{pas} une autre fonction : c'est la \emph{même} fonction, simplement écrite autrement. On passe de l'une à l'autre par développement ou identification.
\end{alerte}

\vfill
\begin{center}
\textcolor{slate}{\itshape\small — Fin du premier portrait. Le suspect a un visage. —}
\end{center}

% ============================================================
%  CHAPITRE 2 — L'AFFAIRE DU BALLON
% ============================================================
\chapterheader{Chapitre 2}{L'affaire du ballon}

Un stade vide, à minuit. Un ballon qui décolle, décrit un arc parfait, et retombe exactement au même endroit. Chaque fois. Comme si quelqu'un avait réglé la trajectoire au millimètre.

L'Inspectrice Maëlle Vasseur savait que derrière chaque trajectoire parfaite se cachait un esprit mathématique. Elle sortit son carnet. La scène de crime avait une équation.

\section{Lire une valeur de \texorpdfstring{$f$}{f}}

\begin{methode}[Calculer $f(x_0)$]
Pour calculer l'image d'un nombre $x_0$ par $f(x) = ax^2+bx+c$, on remplace $x$ par $x_0$ :
\[
  f(x_0) = a\,x_0^2 + b\,x_0 + c.
\]
Aucun piège, mais une rigueur absolue sur les signes.
\end{methode}

\begin{rapport}
Pour $f(x) = 2x^2 - 3x + 1$ :
\[
  f(0) = 1, \qquad f(1) = 2 - 3 + 1 = 0, \qquad f(2) = 8 - 6 + 1 = 3.
\]
\end{rapport}

\section{Le raccourci \texorpdfstring{$f(0) = c$}{f(0) = c}}

\begin{indice}[Le raccourci le plus utile]
Pour toute fonction $f(x) = ax^2+bx+c$ :
\[
  \boxed{\;f(0) = c\;}
\]
Autrement dit, $c$ est \textbf{l'ordonnée à l'origine} : la hauteur à laquelle la parabole coupe l'axe des ordonnées.
\end{indice}

\begin{methode}[Réflexe de l'Inspectrice]
Dès qu'un énoncé donne $f(0)$, l'Inspectrice encadre la valeur de $c$ avant même de lire la suite. C'est cadeau.
\end{methode}

\section{Le tir au mur — la scène de crime}

\begin{indice}[La trajectoire du ballon]
Un ballon lancé depuis le sol a une trajectoire de la forme
\[
  h(t) = at^2 + bt
\]
(ici $c = 0$, car à $t=0$, $h(0)=0$).
\begin{itemize}[leftmargin=*]
  \item Le ballon \textbf{retombe au sol} à la distance $d$ : $h(d) = 0$.
  \item Le ballon \textbf{passe au-dessus d'un mur} situé à la distance $x_w$, de hauteur $h_w$ : $h(x_w) \geq h_w$.
\end{itemize}
\end{indice}

\begin{rapport}[Résolution guidée — drame en quatre actes]
On veut $h(t) = at^2+bt$ avec $h(6) = 0$ et $h(2) \geq 3$.

\textbf{Acte 1 —} $h(6) = 36a + 6b = 0$, donc $b = -6a$.

\textbf{Acte 2 —} On remplace $b$ dans $h(2) \geq 3$ :
\[
  4a + 2b \geq 3 \;\Longrightarrow\; 4a + 2(-6a) \geq 3 \;\Longrightarrow\; -8a \geq 3.
\]

\textbf{Acte 3 —} On divise par $-8 < 0$ : \textbf{le sens de l'inégalité change} !
\[
  a \leq -\frac{3}{8}.
\]

\textbf{Acte 4 —} On \emph{choisit} par exemple $a = -\dfrac{1}{2}$. Alors $b = -6 \times (-\frac{1}{2}) = 3$.

\textbf{Vérification —} $h(6) = -18 + 18 = 0$ \checkmark, et $h(2) = -2 + 6 = 4 \geq 3$ \checkmark.
\end{rapport}

\begin{alerte}[Piège n°4 — le sens de l'inégalité]
Quand on divise (ou multiplie) une inégalité par un nombre \textbf{négatif}, le sens de l'inégalité \textbf{change} :
\[
  -8a \geq 3 \;\Longrightarrow\; a \leq -\frac{3}{8}.
\]
Oublier cette règle est l'erreur la plus fréquente. C'est le \emph{meurtre} de la démonstration.
\end{alerte}

\section{Choix ou déduction ?}

\begin{indice}[Une infinité de solutions]
L'inégalité $a \leq -\dfrac{3}{8}$ admet une \textbf{infinité} de valeurs possibles pour $a$. Passer de cette inégalité à une valeur particulière n'est \textbf{pas une déduction logique} : c'est un \textbf{choix}.

On écrit donc :
\begin{quote}
  « \emph{On choisit par exemple} $a = -\dfrac{1}{2}$ »
\end{quote}
et \textbf{non} « donc $a = -\frac{1}{2}$ » ni « par conséquent $a = -\frac{1}{2}$ ».
\end{indice}

\begin{methode}[Connecteurs logiques — les armes de l'Inspectrice]
\begin{itemize}[leftmargin=*]
  \item \textbf{Déduction nécessaire} : « donc », « par conséquent », « d'où », « en conséquence », « on en déduit que ».
  \item \textbf{Choix} : « on choisit », « prenons par exemple », « on peut prendre ».
\end{itemize}
Utiliser « donc » pour introduire un choix est une \textbf{faute logique}. Une fausse piste.
\end{methode}

\vfill
\begin{center}
\textcolor{slate}{\itshape\small — Le ballon est retombé. Mais un autre crime se prépare. —}
\end{center}

% ============================================================
%  CHAPITRE 3 — LE SCEAU DES ANCIENS
% ============================================================
\chapterheader{Chapitre 3}{Le sceau des anciens}

La deuxième scène était un musée. Une vitrine brisée. Trois chiffres gravés dans la pierre. Sous la vitrine, une seule phrase : \emph{« Les trois coefficients dorment ici. Réveille-les. »}

L'Inspectrice reconnut la signature du suspect. Cette fois, il fallait \textbf{résoudre un système}.

\section{Systèmes linéaires — la méthode de l'Inspectrice}

\begin{methode}[Deux techniques, deux armes]
Pour résoudre un système de deux équations à deux inconnues :
\begin{enumerate}[leftmargin=*]
  \item \textbf{Par substitution} : on exprime une inconnue en fonction de l'autre dans une équation, puis on remplace dans la seconde.
  \item \textbf{Par combinaison (soustraction membre à membre)} : on soustrait les deux équations pour \emph{éliminer} une inconnue.
\end{enumerate}
\end{methode}

\section{Le code à trois chiffres — reconstitution}

\begin{rapport}[L'affaire du sceau]
On cherche $f(x) = ax^2+bx+c$ telle que
\[
  f(0) = 1, \qquad f(1) = 4, \qquad f(2) = 11.
\]

\textbf{Étape 1 —} $f(0) = c = 1$.

\textbf{Étape 2 —} $f(1) = a + b + 1 = 4 \;\Rightarrow\; a + b = 3$.

\textbf{Étape 3 —} $f(2) = 4a + 2b + 1 = 11 \;\Rightarrow\; 2a + b = 5$.

\textbf{Étape 4 —} On soustrait membre à membre :
\[
  (2a+b) - (a+b) = 5 - 3 \;\Longrightarrow\; a = 2.
\]
Puis $b = 3 - a = 1$.

\textbf{Conclusion —} $f(x) = 2x^2 + x + 1$.
\end{rapport}

\begin{alerte}[Piège n°5 — le $c$ oublié]
Ne pas oublier le $c$ dans $f(1)$ : $f(1) = a+b+c$, pas $a+b$ tout seul. C'est l'erreur qui fait s'effondrer tout le raisonnement.
\end{alerte}

\begin{methode}[Stratégie de l'Inspectrice]
Toujours commencer par la valeur la plus simple à exploiter. Ici : $f(0) = 1$ donne \emph{gratuitement} $c = 1$. Une fois $c$ connu, on se ramène à un système $2 \times 2$.
\end{methode}

\vfill
\begin{center}
\textcolor{slate}{\itshape\small — Le sceau est brisé. Trois chiffres livrés. —}
\end{center}

% ============================================================
%  CHAPITRE 4 — LE PONT DES JUMEAUX
% ============================================================
\chapterheader{Chapitre 4}{Le pont des jumeaux}

Troisième scène : un vieux pont de pierre. Deux piliers identiques, plantés de part et d'autre. Entre eux, une parabole tracée au sol. Et sur le parapet, gravé : \emph{« Deux points à la même hauteur. Un axe à découvrir. »}

L'Inspectrice sourit. La symétrie, c'était son terrain.

\section{Deux points à la même ordonnée}

\begin{indice}[Propriété fondamentale de symétrie]
Si deux points de la parabole ont la \textbf{même ordonnée} :
\[
  f(x_1) = f(x_2) \quad \text{avec } x_1 \neq x_2,
\]
alors l'\textbf{axe de symétrie} passe exactement au \textbf{milieu} de $x_1$ et $x_2$ :
\[
  \boxed{\;\alpha = \frac{x_1 + x_2}{2}\;}
\]
\end{indice}

\begin{center}
\begin{tikzpicture}
\begin{axis}[
  width=11cm, height=5.5cm,
  axis lines=middle, xlabel=$x$, ylabel=$y$,
  xmin=-2, xmax=6, ymin=-2, ymax=6,
  xtick={-1,...,5}, ytick={-1,...,5},
  grid=both, grid style={gray!20},
]
\addplot[domain=-1.2:5.2, samples=150, very thick, charcoal] {(x-2)^2 - 1};
\addplot[only marks, mark=*, mark size=3pt, slate] coordinates {(-0.5,5.25) (4.5,5.25)};
\addplot[dashed, very thick, ink] coordinates {(2,-2) (2,6)};
\node at (axis cs:-0.7,5.6) {\small $A$};
\node at (axis cs:4.7,5.6) {\small $B$};
\node at (axis cs:2.4,-1.5) {\small $x=\alpha$};
\end{axis}
\end{tikzpicture}
\end{center}

\section{Retrouver \texorpdfstring{$b$}{b} puis \texorpdfstring{$c$}{c}}

\begin{methode}[Marche à suivre — la reconstitution]
Connaissant $a$ et $\alpha$ :
\begin{enumerate}[leftmargin=*]
  \item On utilise $\alpha = -\dfrac{b}{2a}$ pour isoler $b$ : $b = -2a\alpha$.
  \item On remplace $x$ par les coordonnées d'un point connu pour trouver $c$.
\end{enumerate}
\end{methode}

\begin{rapport}[Le pont des jumeaux]
On sait que $f(x) = ax^2+bx+c$ avec $a = 1$ passe par $A(-1\,;\,0)$ et $B(5\,;\,0)$.

\textbf{Étape 1 —} $A$ et $B$ ont la même ordonnée ($0$), donc
\[
  \alpha = \frac{-1 + 5}{2} = 2.
\]

\textbf{Étape 2 —} $\alpha = -\dfrac{b}{2a} = -\dfrac{b}{2} = 2 \;\Longrightarrow\; b = -4$.

\textbf{Étape 3 —} $f(-1) = 0$ donne $1 - b + c = 0 \Rightarrow 1 + 4 + c = 0 \Rightarrow c = -5$.

\textbf{Conclusion —} $f(x) = x^2 - 4x - 5$.
\end{rapport}

\begin{methode}[Cas particulier $a = 1$]
Si $a = 1$, la relation $\alpha = -\dfrac{b}{2a}$ devient simplement
\[
  \alpha = -\frac{b}{2} \quad\Longleftrightarrow\quad b = -2\alpha.
\]
Un raccourci précieux quand le suspect a un coefficient dominant unitaire.
\end{methode}

\vfill
\begin{center}
\textcolor{slate}{\itshape\small — Le pont tient. L'axe est trouvé. —}
\end{center}

% ============================================================
%  CHAPITRE 5 — LE SOURIRE EMPOISONNÉ
% ============================================================
\chapterheader{Chapitre 5}{Le sourire empoisonné}

Quatrième scène. Un salon bourgeois. Sur le mur, une parabole qui \emph{sourit}. En dessous, un seul indice : $f(x) = ax^2$, avec $a \neq 0$, et $f(2) = -4$.

L'Inspectrice comprit immédiatement : le sourire était un mensonge. La parabole ne souriait pas. Elle \emph{mordait}.

\section{Un seul coefficient}

\begin{indice}[Fonction « carré multiplié »]
Lorsque $b = 0$ et $c = 0$, la fonction s'écrit
\[
  f(x) = ax^2 \quad (a \neq 0).
\]
Sa parabole a pour sommet l'origine $O(0\,;\,0)$ et pour axe de symétrie l'axe des ordonnées ($x = 0$).
\end{indice}

\section{Trouver \texorpdfstring{$a$}{a} à partir d'une valeur}

\begin{methode}[La méthode de l'Inspectrice]
Si l'énoncé donne $f(x_0) = y_0$ avec $x_0 \neq 0$, alors
\[
  a\,x_0^2 = y_0 \;\Longrightarrow\; a = \frac{y_0}{x_0^2}.
\]
\end{methode}

\begin{rapport}
On sait que $f(x) = ax^2$ et $f(3) = -9$. Alors
\[
  a \times 9 = -9 \;\Longrightarrow\; a = -1.
\]
Donc $f(x) = -x^2$.
\end{rapport}

\section{Signe de \texorpdfstring{$f(x)$}{f(x)} et orientation}

\begin{indice}[Signe de $ax^2$]
Pour tout $x \neq 0$, on a $x^2 > 0$. Donc $ax^2$ a le \textbf{même signe que $a$} :
\begin{itemize}[leftmargin=*]
  \item si $a > 0$ : $f(x) > 0$ pour $x \neq 0$, parabole vers le haut ;
  \item si $a < 0$ : $f(x) < 0$ pour $x \neq 0$, parabole vers le bas.
\end{itemize}
Et bien sûr, $f(0) = 0$ dans tous les cas.
\end{indice}

\begin{alerte}[Piège n°6 — le signe qui ne change pas]
Pour $f(x) = ax^2$ avec $a \neq 0$, le signe de $f(x)$ est \textbf{constant} sur $\mathbb{R}^*$. Il ne change pas de signe en cours de route. Un sourire, ou une moue, mais pas les deux.
\end{alerte}

\begin{rapport}[Le sourire empoisonné — reconstitution]
Soit $f(x) = ax^2$ avec $a \neq 0$ et $f(2) = -4$.

\textbf{1.} $4a = -4 \Rightarrow a = -1$.

\textbf{2.} Pour $x \neq 0$, $x^2 > 0$ et $a = -1 < 0$, donc $f(x) = -x^2 < 0$.

\textbf{3.} $a < 0$ : la parabole est tournée vers le bas.
\end{rapport}

\vfill
\begin{center}
\textcolor{slate}{\itshape\small — Le sourire était un piège. Le suspect a encore frappé. —}
\end{center}

% ============================================================
%  CHAPITRE 6 — LE PIC DU DRAGON
% ============================================================
\chapterheader{Chapitre 6}{Le pic du dragon}

Cinquième scène. Une montagne. Un pic acéré. Trois balises lumineuses plantées sur la paroi. Et sur la première balise, gravée : \emph{« Connais le sommet, connais la fonction. »}

L'Inspectrice comprit. Cette fois, le suspect avait laissé tomber son masque. Il fallait reconstituer la \textbf{forme canonique}.

\section{Lire \texorpdfstring{$\alpha$}{alpha} et \texorpdfstring{$\beta$}{beta} dans l'énoncé}

\begin{indice}[Forme canonique]
\[
  f(x) = a(x-\alpha)^2 + \beta
\]
Ici $\alpha$ est l'abscisse du sommet, $\beta$ son ordonnée. On a donc directement $f(\alpha) = \beta$.
\end{indice}

\begin{methode}[Reconstituer sans discriminant]
Étant donnés trois renseignements sur $f$ :
\begin{enumerate}[leftmargin=*]
  \item \textbf{Repérer deux valeurs égales} de $f$. Si $f(x_1) = f(x_2)$, alors
        \[
          \alpha = \frac{x_1+x_2}{2}.
        \]
  \item \textbf{Utiliser la valeur de $f$ au sommet.} Si l'énoncé donne $f(\alpha) = V$, alors $\beta = V$.
  \item \textbf{Utiliser une troisième égalité} pour trouver $a$ :
        \[
          a(x_1-\alpha)^2 + \beta = f(x_1) \;\Longrightarrow\; a = \frac{f(x_1) - \beta}{(x_1-\alpha)^2}.
        \]
\end{enumerate}
\end{methode}

\begin{alerte}[Piège n°7 — le discriminant interdit]
Dans cette énigme, on n'utilise \textbf{pas} le discriminant $\Delta = b^2 - 4ac$. Toute la résolution repose sur la \textbf{symétrie} et la \textbf{substitution}. Le suspect déteste qu'on triche.
\end{alerte}

\section{Le pic du dragon — reconstitution finale}

\begin{rapport}
On cherche $f(x) = a(x-\alpha)^2 + \beta$ telle que
\[
  f(1) = 5, \qquad f(2) = 1, \qquad f(3) = 5.
\]

\textbf{Étape 1 —} On remarque $f(1) = f(3) = 5$. Donc l'axe de symétrie est au milieu :
\[
  \alpha = \frac{1+3}{2} = 2.
\]

\textbf{Étape 2 —} On évalue $f(\alpha) = f(2) = 1$. Donc $\beta = 1$.

\textbf{Étape 3 —} On utilise $f(1) = 5$ :
\[
  a(1-2)^2 + 1 = 5 \;\Longrightarrow\; a \times 1 + 1 = 5 \;\Longrightarrow\; a = 4.
\]

\textbf{Conclusion —}
\[
  \boxed{\,f(x) = 4(x-2)^2 + 1\,}
\]
\end{rapport}

\begin{methode}[Vérification]
Développe : $f(x) = 4(x^2 - 4x + 4) + 1 = 4x^2 - 16x + 17$. Tu peux vérifier que $f(1) = 5$, $f(2) = 1$, $f(3) = 5$.
\end{methode}

\vfill
\begin{center}
\textcolor{slate}{\itshape\small — Le dragon est démasqué. Il ne reste qu'un témoin à interroger. —}
\end{center}

% ============================================================
%  CHAPITRE 7 — L'ART DE L'INTERROGATOIRE
% ============================================================
\chapterheader{Chapitre 7}{L'art de l'interrogatoire}

Sixième scène. Une salle d'interrogatoire. Une table. Deux chaises. Face à l'Inspectrice, un suspect silencieux : \emph{le raisonnement}. L'Inspectrice savait que la vérité ne s'obtient pas par la force, mais par la \textbf{rigueur}.

\section{Les connecteurs logiques}

\begin{indice}[Vocabulaire de l'Inspectrice]
\begin{itemize}[leftmargin=*]
  \item \textbf{Déduction nécessaire} : « donc », « par conséquent », « d'où », « en conséquence », « on en déduit que ».
  \item \textbf{Choix} : « on choisit », « prenons par exemple », « on peut prendre ».
  \item \textbf{Conclusion finale} : « Finalement, \dots » ou « On conclut que \dots ».
\end{itemize}
\end{indice}

\section{Choix ou déduction ? — la règle d'or}

\begin{indice}[Règle d'or]
Lorsqu'une inégalité admet \textbf{une infinité de solutions} et qu'on en sélectionne \emph{une seule}, on ne fait pas une déduction : on fait un \textbf{choix}. La phrase doit le signaler explicitement.
\end{indice}

\begin{rapport}[Comparaison]
\textbf{À éviter :} « $-8a \geq 3$, donc $a = -\frac{1}{2}$. » \\
\textbf{Correct :} « $-8a \geq 3$, donc $a \leq -\frac{3}{8}$. On choisit par exemple $a = -\frac{1}{2}$. »
\end{rapport}

\section{Structure type d'une réponse}

\begin{methode}[Modèle de rédaction]
\begin{enumerate}[leftmargin=*]
  \item \textbf{Ce que je sais} : je recopie les données et je les traduis en équations.
  \item \textbf{Ce que j'en déduis} : j'enchaîne les étapes avec des connecteurs corrects.
  \item \textbf{Ce que je choisis} (s'il y a lieu) : je signale explicitement tout choix.
  \item \textbf{Conclusion encadrée} : je récapitule clairement le résultat.
\end{enumerate}
\end{methode}

\begin{alerte}[Piège n°8 — le connecteur menteur]
Un « donc » là où il n'y a qu'un choix, c'est un \emph{faux témoignage}. L'Inspectrice ne le tolère pas.
\end{alerte}

\vfill
\begin{center}
\textcolor{slate}{\itshape\small — Le raisonnement a parlé. Le dossier se referme. —}
\end{center}

% ============================================================
%  CHAPITRE 8 — LE CARNET DE L'INSPECTRICE
% ============================================================
\chapterheader{Chapitre 8}{Le carnet de l'Inspectrice}

L'aube pointait. L'Inspectrice Maëlle Vasseur ouvrit son carnet et relut ses notes. Six affaires. Six signatures. Un seul suspect.

\section{Formulaire — tout ce qu'il faut retenir}

\begin{indice}[Le carnet de l'Inspectrice]
\begin{align*}
  &\text{Forme développée :} && f(x) = ax^2+bx+c \quad (a \neq 0) \\[2pt]
  &\text{Forme canonique :} && f(x) = a(x-\alpha)^2+\beta \\[2pt]
  &\text{Abscisse du sommet :} && \alpha = -\frac{b}{2a} \\[2pt]
  &\text{Ordonnée du sommet :} && \beta = f(\alpha) \\[2pt]
  &\text{Axe de symétrie :} && x = \alpha \\[2pt]
  &\text{Ordonnée à l'origine :} && f(0) = c \\[2pt]
  &\text{Points de même hauteur :} && \alpha = \frac{x_1+x_2}{2} \text{ si } f(x_1)=f(x_2) \\[2pt]
  &\text{Lien entre $\alpha$ et $b$ :} && b = -2a\alpha \\[2pt]
  &\text{Orientation :} && a>0 \Rightarrow \text{vers le haut} \quad ; \quad a<0 \Rightarrow \text{vers le bas}
\end{align*}
\end{indice}

\section{Les huit réflexes de l'Inspectrice}

\begin{methode}[Mémo]
\begin{enumerate}[leftmargin=*]
  \item Je vois $f(0)$ ? $\Rightarrow$ je note $c$ immédiatement.
  \item Je vois deux valeurs égales de $f$ ? $\Rightarrow$ axe de symétrie au milieu.
  \item Je cherche $\alpha$ ? $\Rightarrow$ $\alpha = -b/(2a)$.
  \item Je divise une inégalité par un négatif ? $\Rightarrow$ je change le sens.
  \item J'ai une infinité de solutions ? $\Rightarrow$ je \emph{choisis} et je l'écris.
  \item On me donne un sommet ? $\Rightarrow$ j'utilise la forme canonique.
  \item On me donne $f(x_0)$ ? $\Rightarrow$ je substitue, sans oublier $c$.
  \item On me dit « sans discriminant » ? $\Rightarrow$ je passe par la symétrie.
\end{enumerate}
\end{methode}

\section{Exercices récapitulatifs — les affaires classées}

\begin{affaire}[Affaire n°1 — Le ballon]
Un ballon a une trajectoire $h(t) = at^2+bt$. Il retombe au sol à $8$ m et passe au-dessus d'un mur situé à $3$ m, de hauteur $4$ m.
\begin{enumerate}[leftmargin=*]
  \item Traduire les conditions en équations/inéquations.
  \item Exprimer $b$ en fonction de $a$.
  \item En déduire un encadrement de $a$.
  \item Choisir une valeur de $a$ et conclure.
\end{enumerate}
\end{affaire}

\begin{revelation}
\textbf{1.} $h(8) = 0$ et $h(3) \geq 4$. \\
\textbf{2.} $64a + 8b = 0 \Rightarrow b = -8a$. \\
\textbf{3.} $h(3) = 9a + 3(-8a) = -15a \geq 4 \Rightarrow a \leq -\frac{4}{15}$. \\
\textbf{4.} On choisit $a = -\frac{1}{2}$, alors $b = 4$. Vérification : $h(8) = -32+32 = 0$ \checkmark, $h(3) = -4{,}5+12 = 7{,}5 \geq 4$ \checkmark.
\end{revelation}

\begin{affaire}[Affaire n°2 — Le code]
Trouver $f(x) = ax^2+bx+c$ telle que $f(0) = 2$, $f(1) = 3$, $f(2) = 8$.
\end{affaire}

\begin{revelation}
$c = 2$. Puis $a+b = 1$ et $2a+b = 3$. Par soustraction : $a = 2$, puis $b = -1$. Donc $f(x) = 2x^2 - x + 2$.
\end{revelation}

\begin{affaire}[Affaire n°3 — Le sourire]
Soit $f(x) = ax^2$ avec $a \neq 0$ et $f(4) = -8$.
\begin{enumerate}[leftmargin=*]
  \item Déterminer $a$.
  \item Donner le signe de $f(x)$ pour $x \neq 0$.
  \item Donner l'orientation de la parabole.
\end{enumerate}
\end{affaire}

\begin{revelation}
\textbf{1.} $16a = -8 \Rightarrow a = -\frac{1}{2}$. \\
\textbf{2.} $a<0$ et $x^2 >0$ pour $x \neq 0$, donc $f(x) < 0$. \\
\textbf{3.} $a<0$ : parabole vers le bas.
\end{revelation}

\begin{affaire}[Affaire n°4 — Les jumeaux]
Une parabole $f(x) = x^2 + bx + c$ passe par $A(-2\,;\,0)$ et $B(6\,;\,0)$.
\begin{enumerate}[leftmargin=*]
  \item Déterminer $\alpha$.
  \item En déduire $b$.
  \item Calculer $c$.
\end{enumerate}
\end{affaire}

\begin{revelation}
\textbf{1.} $\alpha = \dfrac{-2+6}{2} = 2$. \\
\textbf{2.} $2 = -\dfrac{b}{2} \Rightarrow b = -4$. \\
\textbf{3.} $f(-2) = 4 + 8 + c = 0 \Rightarrow c = -12$.
\end{revelation}

\begin{affaire}[Affaire n°5 — Le dragon]
On cherche $f(x) = a(x-\alpha)^2+\beta$ telle que $f(0) = 5$, $f(2) = 1$, $f(4) = 5$.
\end{affaire}

\begin{revelation}
$f(0) = f(4) = 5 \Rightarrow \alpha = 2$. Puis $f(2) = 1 \Rightarrow \beta = 1$. Enfin $f(0) = a(0-2)^2 + 1 = 5 \Rightarrow 4a = 4 \Rightarrow a = 1$. Donc $f(x) = (x-2)^2 + 1$.
\end{revelation}

\vfill
\begin{center}
\textcolor{slate}{\itshape\small — Toutes les affaires sont classées. —}
\end{center}

% ============================================================
%  ÉPILOGUE
% ============================================================
\chapterheader{Épilogue}{La dernière pièce}

L'aube se levait sur la ville. L'Inspectrice Maëlle Vasseur referma le dernier dossier. Six affaires. Six paraboles. Un seul coupable, finalement : \emph{le second degré lui-même}.

Mais l'Inspectrice ne souriait pas. Elle savait que derrière chaque parabole se cachait une nouvelle énigme. Et que la nuit reviendrait.

Elle rangea son carnet. Sur la couverture, une phrase, écrite à la main :
\[
  \boxed{\;\text{« Connais le sommet, connais la fonction. »}\;}
\]

\vspace{1cm}

\begin{indice}[Ce que vous avez appris cette nuit]
\begin{itemize}[leftmargin=*]
  \item la forme développée $ax^2+bx+c$ et la forme canonique $a(x-\alpha)^2+\beta$ ;
  \item le rôle du signe de $a$ (orientation) ;
  \item le sommet $(\alpha\,;\,\beta)$ et l'axe de symétrie $x=\alpha$ ;
  \item l'exploitation de $f(0) = c$ et des systèmes linéaires ;
  \item la symétrie des points de même ordonnée ;
  \item la rédaction rigoureuse d'un raisonnement (choix \emph{vs} déduction).
\end{itemize}
\end{indice}

\vspace{1cm}

\begin{center}
\begin{tcolorbox}[
  enhanced, colback=ink, colframe=ink,
  interior style={top color=ink, bottom color=charcoal, shading angle=90},
  boxrule=0pt, arc=2mm, width=0.8\textwidth,
  left=18pt, right=18pt, top=14pt, bottom=14pt,
]
\centering
{\color{silver}\itshape « Toute parabole cache un crime.\\[2pt]
Toute déduction cache un choix.\\[2pt]
Toute conclusion cache une nuit blanche. »}\\[8pt]
{\color{gray}\small — Inspectrice Maëlle Vasseur, 23h47}
\end{tcolorbox}
\end{center}

\vfill

\begin{center}
\textcolor{slate}{\small\itshape Fin du dossier.\\[2pt]
\textbf{L'Odyssée du Second Degré} — Yann Merdy, 2026.\\[2pt]
Document pédagogique — Licence non commerciale.}
\end{center}

\end{document}